\ea\label{ex:text14-1}
A takinima eve a peha taba to tabotaboo vai niara nomanai en Teapu.\\
\gll a	takinima	eve	a	peha	taba	to	tabo{$\sim$}taboo	vai	ni=ara	nomana=i	en	teapu\\
     \textsc{art}2.\textsc{sg}	woven.bamboo	3\textsc{sg}.\textsc{pron}	\textsc{art}2.\textsc{sg}	\textsc{indef}2.\textsc{sg}	thing	\textsc{rel}	\textsc{red}{$\sim$}make.wall	now	\textsc{appl}=1\textsc{pl}.\textsc{in}.\textsc{pron}	today=\textsc{dem}1	\textsc{dem}3	Teop.region\\
\glt `Woven bamboo is a thing that we use today for making walls here in the Teop region.'
\z

\ea\label{ex:text14-2}
Eve he na paku raara bona koopu.\\
\gll eve	he	na	paku	ra=ara	bona	koopu\\
     3\textsc{sg}.\textsc{pron}	\textsc{advers}	\textsc{tm}2	make	1\textsc{pl}.\textsc{in}.\textsc{ipfv}=1\textsc{pl}.\textsc{in}.\textsc{pron}	\textsc{obj}.\textsc{art}3.\textsc{pl}	koopu.bamboo\\
\glt `We make it from \textit{koopu} bamboos.'
\z

\ea\label{ex:text14-3}
A kiu nae he na paku vonaen:\\
\gll a	kiu	nae	he	na	paku	vo=na=en\\
     \textsc{art}2.\textsc{sg}	work	3\textsc{sg}.\textsc{poss}.\textsc{pron}	\textsc{advers}	\textsc{tm}2	do	like=3\textsc{sg}.\textsc{ipfv}=\textsc{dem}3\\
\glt `The work on it is done as follows::'
\z

\ea\label{ex:text14-4}
Eara na tagihu rara rasuu\\
\gll e-ara	na	tagihu	rara	rasuu\\
     \textsc{h}-1\textsc{pl}.\textsc{in}.\textsc{pron}	\textsc{tm}2	go.into.the.jungle	1\textsc{pl}.\textsc{in}.\textsc{ipfv}	jungle\\
\glt `We go into the jungle'
\z

\ea\label{ex:text14-5}
ei kou a koopu sa teitei vira bata haana kasuana.\\
\gll ei~kou	a	koopu	sa	tei{$\sim$}tei	vira	bata	haa=na	kasuana\\
     because	\textsc{art}3.\textsc{pl}	koopu.bamboo	\textsc{neg}1	\textsc{red}{$\sim$}be.there	much	\textsc{cont}	\textsc{neg}2=3\textsc{sg}.\textsc{ipfv}	coast\\
\glt `because there are not many \textit{koopu} bamboos along the coast.'
\z

\ea\label{ex:text14-6}
A koopu vai a mataa ae bebeahu na teitei nana rasuu.\\
\gll a	koopu	vai	a	mataa	ae	bebeahu	na	tei{$\sim$}tei	nana	rasuu\\
     \textsc{art}3.\textsc{pl}	koopu.bamboo	\textsc{dem}6	\textsc{art}3.\textsc{pl}	good	and1	long	\textsc{tm}2	\textsc{red}{$\sim$}be.there	3\textsc{sg}.\textsc{ipfv}	jungle\\
\glt `The good and long \textit{koopu} bamboos grow in the jungle.'
\z

\ea\label{ex:text14-7}
Beara tagihu ei to sun nae a koopu,\\
\gll beara	tagihu	ei	to	sun	nae	a	koopu\\
     when.1\textsc{pl}.\textsc{in}.\textsc{pron}	go.into.the.jungle	\textsc{dem}2	\textsc{rel}	stand	3\textsc{sg}.\textsc{ipfv}	\textsc{art}3.\textsc{pl}	koopu.bamboo\\
\glt `When we go into the jungle where the \textit{koopu} bamboos stand,'
\z

\ea\label{ex:text14-8}
eara toro siisira rakaha roho a koopu eara repaa mosi mana.\\
\gll e-ara	toro	siisira	rakaha	roho	a	koopu	e-ara	re=paa	mosi	mana\\
     \textsc{h}-1\textsc{pl}.\textsc{in}.\textsc{pron}	\textsc{tm}7	examine	intensively	before	\textsc{art}3.\textsc{pl}	koopu.bamboo	\textsc{h}-1\textsc{pl}.\textsc{in}.\textsc{pron}	\textsc{consec}=\textsc{tm}3	cut	eventually\\
\glt `we must intensively examine the \textit{koopu} bamboos, before we eventually cut them.'
\z

\ea\label{ex:text14-9}
O koopu toro mosi ara toro manoto bara toro mahaka rakaha mee.\\
\gll o	koopu	toro	mosi	ara	toro	manoto	bara	toro	mahaka	rakaha	mee\\
     \textsc{art}3.\textsc{sg}	koopu.bamboo	\textsc{tm}7	cut	1\textsc{pl}.\textsc{in}.\textsc{pron}	\textsc{tm}7	mature	and2	\textsc{tm}7	dry	very	also\\
\glt `The \textit{koopu} bamboo that we should cut must be mature and also must be very dry.'
\z

\ea\label{ex:text14-10}
Eara sa antee haara tea mosi o koopu vai o kobus\\
\gll e-ara	sa	antee	haa=ra	tea	mosi	o	koopu	vai	o	kobus\\
     \textsc{h}-1\textsc{pl}.\textsc{in}.\textsc{pron}	\textsc{neg}1	can	\textsc{neg}2=1\textsc{pl}.\textsc{in}.\textsc{ipfv}	\textsc{compl}1	cut	\textsc{art}3.\textsc{sg}	koopu.bamboo	\textsc{dem}6	\textsc{art}3.\textsc{sg}	moist\\
\glt `We cannot cut the moist \textit{koopu} bamboo,'
\z

\ea\label{ex:text14-11}
tabae na antee nana tea koava ge mahamahan\\
\gll tabae	na	antee	nana	tea	koava	ge	maha{$\sim$}mahan\\
     because	\textsc{tm}2	can	3\textsc{sg}.\textsc{ipfv}	\textsc{compl}1	rot	or	\textsc{red}{$\sim$}space\\
\glt `because it can rot or leave spaces,'
\z

\ea\label{ex:text14-12}
beara takinima nie\\
\gll be=ara	takinima	ni=e\\
     when=1\textsc{pl}.\textsc{in}.\textsc{pron}	weave.bamboo.wall	\textsc{appl}=3\textsc{sg}.\textsc{pron}\\
\glt `when we weave the bamboo wall with it,'
\z

\ea\label{ex:text14-13}
ei be raka bata maa a koopu.\\
\gll ei	be	raka	bata	maa	a	koopu\\
     \textsc{dem}2	when	dry	\textsc{cont}	hither	\textsc{art}3.\textsc{pl}	koopu.bamboo\\
\glt `that is when the \textit{koopu} bamboos dry up.'
\z

\ea\label{ex:text14-14}
Beara siisira vakavara, eara repaa mosi.\\
\gll beara	siisira	va-kavara	e-ara	re=paa	mosi\\
     when.1\textsc{pl}.\textsc{in}.\textsc{pron}	examine	\textsc{advr}-complete	\textsc{h}-1\textsc{pl}.\textsc{in}.\textsc{pron}	\textsc{consec}=\textsc{tm}3	cut\\
\glt `When we have finished examining (them), we cut (them).'
\z

\ea\label{ex:text14-15}
Are mosi vaantee a inu, eara repaa paku amaa dee.\\
\gll a=re	mosi	va-antee	a	inu	e-ara	re=paa	paku	a=maa	dee\\
     1\textsc{pl}.\textsc{in}.\textsc{pron}=\textsc{consec}	cut	\textsc{caus}-enough	\textsc{art}2.\textsc{sg}	house	\textsc{h}-1\textsc{pl}.\textsc{in}.\textsc{pron}	\textsc{consec}=\textsc{tm}3	make	\textsc{art}2.\textsc{sg}=\textsc{plm}	load\\
\glt `Then we cut as many as are needed for the house, and then we make the loads (to carry the bamboos home).'
\z

\ea\label{ex:text14-16}
A peha dee kou na antee batana bono peho saavun koopu.\\
\gll a	peha	dee	kou	na	antee	bata=na	bono	peho	saavun	koopu\\
     \textsc{art}2.\textsc{sg}	one2.\textsc{sg}	load	\textsc{emph}2	\textsc{tm}2	consist.of	\textsc{cont}=3\textsc{sg}.\textsc{ipfv}	\textsc{obj}.\textsc{art}3.\textsc{sg}	one3.\textsc{sg}	ten	koopu.bamboo\\
\glt `One load consists of ten \textit{koopu} bamboos.'
\z

\ea\label{ex:text14-17}
Kavara vai, eara dee voosu vo vaan.\\
\gll kavara	vai	e-ara	dee	voosu	vo	vaan\\
     finished	now	\textsc{h}-1\textsc{pl}.\textsc{in}.\textsc{pron}	carry	go.home	\textsc{goal}	village\\
\glt `After that we carry (them) home to the village.'
\z

\ea\label{ex:text14-18}
Ei vaan to kahi paku pete raara o hum vaa tea kopakopaa a koopu.\\
\gll ei	vaan	to	kahi	paku	pete	ra=ara	o	hum	vaa	tea	kopa{$\sim$}kopaa	a	koopu\\
     \textsc{dem}2	village	\textsc{rel}	\textsc{tm}1	do	subsequently	1\textsc{pl}.\textsc{in}.\textsc{ipfv}=1\textsc{pl}.\textsc{in}.\textsc{pron}	\textsc{art}3.\textsc{sg}	part	\textsc{lnk}	\textsc{compl}1	\textsc{red}{$\sim$}cut.lengthwise	\textsc{art}3.\textsc{pl}	koopu.bamboo\\
\glt `It is here in the village where we will do the next part o work, the lengthwise cutting of the \textit{koopu} bamboos.'
\z

\ea\label{ex:text14-19}
Beara kopakopaa vakavara, eara repaa paku pete o hum vaa tea tasu.\\
\gll beara	kopa{$\sim$}kopaa	va-kavara	e-ara	re=paa	paku	pete	o	hum	vaa	tea	tasu\\
     when.1\textsc{pl}.\textsc{in}.\textsc{pron}	\textsc{red}{$\sim$}cut.lengthwise	\textsc{advr}-complete	\textsc{h}-1\textsc{pl}.\textsc{in}.\textsc{pron}	\textsc{consec}=\textsc{tm}3	do	subsequently	\textsc{art}3.\textsc{sg}	part	\textsc{lnk}	\textsc{compl}1	beat\\
\glt `When we have cut (them) lengthwise, we do the next part of (work), the beating (of the bamboo strips).'\footnote{lit. `..., then we subsequently do the part of beating.'}
\z

\ea\label{ex:text14-20}
A tasu ei beara tasu sirisiri bata rara a koopu, ore tamee teo vuri beara takitakinima batara.\\
\gll a	tasu	ei	be=ara	tasu	siri{$\sim$}siri	bata	rara	a	koopu	o=re	tamee	te=o	vuri	be=ara	taki{$\sim$}takinima	bata=ra\\
     \textsc{art}2.\textsc{sg}	beat	\textsc{dem}2	\textsc{compl}2=1\textsc{pl}.\textsc{in}.\textsc{pron}	beat	\textsc{red}{$\sim$}split	\textsc{cont}	1\textsc{pl}.\textsc{in}.\textsc{ipfv}	\textsc{art}3.\textsc{pl}	koopu.bamboo	3.\textsc{pron}=\textsc{consec}	flexible	\textsc{prep}=\textsc{art}3.\textsc{sg}	time	when=1\textsc{pl}.\textsc{in}.\textsc{pron}	\textsc{red}{$\sim$}weave.bamboo	\textsc{cont}=1\textsc{pl}.\textsc{in}.\textsc{ipfv}\\
\glt `This beating (means) that we beat (and) split the \textit{koopu} bamboo strips, so that they are flexible when we are weaving the bamboo wall.'
\z

\ea\label{ex:text14-21}
Tea tasu he eara na tasu rara a koopu bona toraara.\\
\gll te=a	tasu	he	e-ara	na	tasu	rara	a	koopu	bona	toraara\\
     \textsc{prep}=\textsc{art}2.\textsc{sg}	beat	\textsc{advers}	\textsc{h}-1\textsc{pl}.\textsc{in}.\textsc{pron}	\textsc{tm}2	beat	1\textsc{pl}.\textsc{in}.\textsc{ipfv}	\textsc{art}3.\textsc{pl}	koopu.bamboo	\textsc{obj}.\textsc{art}2.\textsc{sg}	axe\\
\glt `As for the beating, we beat the \textit{koopu} bamboo strips with an axe.'
\z

\ea\label{ex:text14-22}
Be kavara a kiu vaa tea tasu, eara repaa taatagi o komano koopu -\\
\gll be	kavara	a	kiu	vaa	tea	tasu	e-ara	re=paa	taatagi	o	koma-n=o	koopu\\
     when	finished	\textsc{art}2.\textsc{sg}	work	\textsc{lnk}	\textsc{compl}1	beat	\textsc{h}-1\textsc{pl}.\textsc{in}.\textsc{pron}	\textsc{consec}=\textsc{tm}3	prepare	\textsc{art}2.\textsc{pl}	inside-3\textsc{sg}.\textsc{poss}=\textsc{art}3.\textsc{sg}	koopu.bamboo\\
\glt `When the work of the beating is finished, we prepare the inside of the \textit{koopu} bamboo strips.'
\z

\newpage
\ea\label{ex:text14-23}
ei eara na taraha vatotobin bata koara amaa nava nae.\\
\gll ei	e-ara	na	taraha	va-totobin	bata	koa=ra	a=maa	nava	na=e\\
     \textsc{dem}2	\textsc{h}-1\textsc{pl}.\textsc{in}.\textsc{pron}	\textsc{tm}2	shave	\textsc{caus}-even	\textsc{cont}	just=1\textsc{pl}.\textsc{in}.\textsc{ipfv}	\textsc{art}2.\textsc{sg}=\textsc{plm}	bar	3\textsc{sg}.\textsc{poss}=3\textsc{sg}.\textsc{pron}\\
\glt `this (means that) we shave (and) make its bars even,'
\z

\ea\label{ex:text14-24}
Ei to mene ato vakisi bata ni nana bono koopu teo vuri vaa tea takitakinima.\\
\gll ei	to	mene	ato	vakisi	bata	ni	nana	bono	koopu	te=o	vuri	vaa	tea	taki{$\sim$}takinima\\
     \textsc{dem}2	\textsc{rel}	\textsc{tm}6	hold	still	\textsc{cont}	\textsc{appl}	3\textsc{sg}.\textsc{ipfv}	\textsc{obj}.\textsc{art}3.\textsc{sg}	koopu.bamboo	\textsc{prep}=\textsc{art}3.\textsc{sg}	time	\textsc{lnk}	\textsc{compl}1	\textsc{red}{$\sim$}weave.bamboo.wall\\
\glt `This is to prevent a \textit{koopu} bamboo strips from getting stuck (at the knots) when the weaving is done.'\footnote{lit. `This is what otherwise would hold the bamboo at the time of weaving.'}
\z

\ea\label{ex:text14-25}
Beara taatagi vakavara a koopu, eara nao vai tea kiu vaa tea takitakinima.\\
\gll beara	taatagi	va-kavara	a	koopu	e-ara	nao	vai	te=a	kiu	vaa	tea	taki{$\sim$}takinima\\
     when.1\textsc{pl}.\textsc{in}.\textsc{pron}	prepare	\textsc{advr}-complete	\textsc{art}3.\textsc{pl}	koopu.bamboo	\textsc{h}-1\textsc{pl}.\textsc{in}.\textsc{pron}	go	now	\textsc{prep}=\textsc{art}2.\textsc{sg}	work	\textsc{lnk}	\textsc{compl}1	\textsc{red}{$\sim$}weave.bamboo.wall\\
\glt `When we have prepared the \textit{koopu} bamboo strips, we proceed with the work of the bamboo-weaving.'
\z

\ea\label{ex:text14-26}
Tea takitakinima he na tei me nana bono buaku o kaku hanana vaa tea paku takinima.\\
\gll te=a	taki{$\sim$}takinima	he	na	tei	me	nana	bono	buaku	o	kaku	hanana	vaa	tea	paku	takinima\\
     \textsc{prep}=\textsc{art}2.\textsc{sg}	\textsc{red}{$\sim$}weave.bamboo	\textsc{advers}	\textsc{tm}2	stay	with	3\textsc{sg}.\textsc{ipfv}	\textsc{obj}.\textsc{art}3.\textsc{sg}	two	\textsc{art}3.\textsc{sg}	kind	pattern	\textsc{lnk}	\textsc{compl}1	make	bamboo.wall\\
\glt `As for the bamboo-weaving, there are two kinds of pattern of making woven bamboo walls.'
\z

\newpage
\ea\label{ex:text14-27}
O peho o vavaasun, o meho he o vavaagoroho.\\
\gll o	peho	o	va{$\sim$}vaa-sun	o	meho	he	o	va{$\sim$}vaa-goroho\\
     \textsc{art}3.\textsc{sg}	one3.\textsc{sg}	\textsc{art}3.\textsc{sg}	\textsc{red}{$\sim$}\textsc{caus}-stand	\textsc{art}3.\textsc{sg}	other3.\textsc{sg}	\textsc{advers}	\textsc{art}3.\textsc{sg}	\textsc{red}{$\sim$}\textsc{caus}-sleep\\
\glt `One is the make-stand-up pattern, the other one the make-asleep pattern.'
\z

\ea\label{ex:text14-28}
Teo vavaasun, eara na ipinoo rara a sun hiava bara a tatabonata na inu.\\
\gll te=o	va{$\sim$}vaa-sun	e-ara	na	ipinoo	rara	a	sun	hiava	bara	a	tatabonata	n=a	inu\\
     \textsc{prep}=\textsc{art}3.\textsc{sg}	\textsc{red}{$\sim$}\textsc{caus}-stand	\textsc{h}-1\textsc{pl}.\textsc{in}.\textsc{pron}	\textsc{tm}2	measure	1\textsc{pl}.\textsc{in}.\textsc{ipfv}	\textsc{art}2.\textsc{sg}	stand	go.up	and2	\textsc{art}2.\textsc{sg}	width	3\textsc{sg}.\textsc{poss}=\textsc{art}2.\textsc{sg}	house\\
\glt `As for the make-stand-up pattern, we measure the height and the width of the house.'
\z

\ea\label{ex:text14-29}
Eara repaa kurukurus a koopu vai to paa tasu vaevuru ara.\\
\gll e-ara	re=paa	kuru{$\sim$}kurus	a	koopu	vai	to	paa	tasu	vaevuru	ara\\
     \textsc{h}-1\textsc{pl}.\textsc{in}.\textsc{pron}	\textsc{consec}=\textsc{tm}3	\textsc{red}{$\sim$}cut	\textsc{art}3.\textsc{pl}	koopu.bamboo	\textsc{dem}6	\textsc{rel}	\textsc{tm}3	beat	already	1\textsc{pl}.\textsc{in}.\textsc{pron}\\
\glt `Then we cut the \textit{koopu} bamboo strips that we have already beaten (according to the measurements).'
\z

\ea\label{ex:text14-30}
Be kavara o hum bona, eara repaa vavaasun nao a koopu, ore no suguna komana teo vagegeana na taboo.\\
\gll be	kavara	o	hum	bona	e-ara	re=paa	va{$\sim$}vaa-sun	nao	a	koopu	o=re	no	suguna	komana	te=o	vagegeana	n=a	taboo\\
     when	finished	\textsc{art}3.\textsc{sg}	part	\textsc{ana}.\textsc{sg}	\textsc{h}-1\textsc{pl}.\textsc{in}.\textsc{pron}	\textsc{consec}=\textsc{tm}3	\textsc{red}{$\sim$}\textsc{caus}-stand	thither	\textsc{art}3.\textsc{pl}	koopu.bamboo	3.\textsc{pron}=\textsc{consec}	go	arrive	altogether	\textsc{prep}=\textsc{art}2.\textsc{pl}	end	3\textsc{sg}.\textsc{poss}=\textsc{art}2.\textsc{sg}	wall\\
\glt `When this part is finished, we make the \textit{koopu} bamboo strips stand up, until the ends of the wall are reached.'
\z

\ea\label{ex:text14-31}
Eara repaa taneo tea hogohogo ni nao a koopu vaantee tea tatabonata na taboo.\\
\gll e-ara	re=paa	taneo	tea	hogo{$\sim$}hogo	ni	nao	a	koopu	va-antee	te=a	tatabonata	n=a	taboo\\
     \textsc{h}-1\textsc{pl}.\textsc{in}.\textsc{pron}	\textsc{consec}=\textsc{tm}3	start	\textsc{compl}1	\textsc{red}{$\sim$}push	\textsc{appl}	thither	\textsc{art}3.\textsc{pl}	koopu.bamboo	\textsc{caus}-enough	\textsc{prep}=\textsc{art}2.\textsc{sg}	width	3\textsc{sg}.\textsc{poss}=\textsc{art}2.\textsc{sg}	wall\\
\glt `Then we start pushing the (horizontal) \textit{koopu} bamboo strips (in between the vertical ones), until it is enough for the width of the wall.'
\z

\ea\label{ex:text14-32}
Beara hogo, eara na gae rara a ta kukan, eara repaa hogo tavusu nie.\\
\gll beara	hogo	e-ara	na	gae	rara	a	ta	kukan	e-ara	re=paa	hogo	tavusu	ni=e\\
     when.1\textsc{pl}.\textsc{in}.\textsc{pron}	push	\textsc{h}-1\textsc{pl}.\textsc{in}.\textsc{pron}	\textsc{tm}2	lift	1\textsc{pl}.\textsc{in}.\textsc{ipfv}	\textsc{art}2.\textsc{sg}	piece	three	\textsc{h}-1\textsc{pl}.\textsc{in}.\textsc{pron}	\textsc{consec}=\textsc{tm}3	push	come.out	\textsc{appl}=3\textsc{sg}.\textsc{pron}\\
\glt `When we push them, we lift three (vertical) strips, then we push them (the horizontal ones) through.'
\z

\ea\label{ex:text14-33}
Eara repaa taono a ta meha a ta kukan, eara repaa gae vaha bana a ta meha ta kukan.\\
\gll e-ara	re=paa	taono	a	ta	meha	a	ta	kukan	e-ara	re=paa	gae	vahabana	a	ta	meha	ta	kukan\\
     \textsc{h}-1\textsc{pl}.\textsc{in}.\textsc{pron}	\textsc{consec}=\textsc{tm}3	put.on.top	\textsc{art}2.\textsc{sg}	piece	other2.\textsc{sg}	\textsc{art}2.\textsc{sg}	piece	three	\textsc{h}-1\textsc{pl}.\textsc{in}.\textsc{pron}	\textsc{consec}=\textsc{tm}3	lift	again	\textsc{art}2.\textsc{sg}	piece	other2.\textsc{sg}	piece	three\\
\glt `Then we put the three other strips on top, and then we lift another three strips again.'
\z

\ea\label{ex:text14-34}
Eara paku bata koa vo nao raen ore poosu komana a taboo.\\
\gll e-ara	paku	bata	koa	vo	nao	ra=en	o=re	poosu	komana	a	taboo\\
     \textsc{h}-1\textsc{pl}.\textsc{in}.\textsc{pron}	do	\textsc{cont}	just	like	thither	1\textsc{pl}.\textsc{in}.\textsc{ipfv}=\textsc{dem}3	3.\textsc{pron}=\textsc{consec}	be.closed	altogether	\textsc{art}2.\textsc{sg}	wall\\
\glt `We just do it like this until the wall is completely closed.'
\z

\ea\label{ex:text14-35}
Teo vavaagoroho he na vaapeha koana, amaa hogo ae a gae kukan.\\
\gll te=o	va{$\sim$}vaa-goroho	he	na	vaa-peha	koa=na	a=maa	hogo	ae	a	gae	kukan\\
     \textsc{prep}=\textsc{art}3.\textsc{sg}	\textsc{red}{$\sim$}\textsc{caus}-sleep	\textsc{advers}	\textsc{tm}2	\textsc{caus}-one	just=3\textsc{sg}.\textsc{ipfv}	\textsc{art}2.\textsc{sg}=\textsc{plm}	push	and1	\textsc{art}2.\textsc{sg}	lift	three\\
\glt `With the make-asleep pattern it is just the same, some pushing and the lifting of three (strips).'
\z

\ea\label{ex:text14-36}
Evehee ei kou eara sa paku haara to ipinoo.\\
\gll evehee	ei	kou	e-ara	sa	paku	haa=ra	to	ipinoo\\
     but	\textsc{dem}2	\textsc{emph}2	\textsc{h}-1\textsc{pl}.\textsc{in}.\textsc{pron}	\textsc{neg}1	do	\textsc{neg}2=1\textsc{pl}.\textsc{in}.\textsc{ipfv}	\textsc{nspec}3.\textsc{sg}	measure\\
\glt `But now we do not do any measurement.'
\z

\ea\label{ex:text14-37}
Eara na vaagoroho rara a ta kukan ge a ta taonim,\\
\gll e-ara	na	vaa-goroho	rara	a	ta	kukan	ge	a	ta	taonim\\
     \textsc{h}-1\textsc{pl}.\textsc{in}.\textsc{pron}	\textsc{tm}2	\textsc{caus}-sleep	1\textsc{pl}.\textsc{in}.\textsc{ipfv}	\textsc{art}2.\textsc{sg}	piece	three	or	\textsc{art}2.\textsc{sg}	piece	five\\
\glt `We lay three or five strips down,'
\z

\ea\label{ex:text14-38}
eara repaa taneo kou tea hogo ae gae, ore suguna komana bea taboo poosu vakavara.\\
\gll e-ara	re=paa	taneo	kou	tea	hogo	ae	gae	o=re	suguna	komana	be=a	taboo	poosu	va-kavara\\
     \textsc{h}-1\textsc{pl}.\textsc{in}.\textsc{pron}	\textsc{consec}=\textsc{tm}3	start	\textsc{emph}2	\textsc{compl}1	push	and1	lift	3.\textsc{pron}=\textsc{consec}	arrive	altogether	\textsc{compl}2=\textsc{art}2.\textsc{sg}	wall	be.closed	\textsc{advr}-complete\\
\glt `then we start pushing and lifting (the strips) until the wall is completely closed.'\footnote{lit. `... then (the time) arrives that the wall is completely closed.'}
\z

\ea\label{ex:text14-39}
Ahiki to ipinoo, eikou eara na ponuu kaakoo bata koara a koopu beara hogo vaantee bata naora a taboo.\\
\gll ahiki	to	ipinoo	eikou	e-ara	na	ponuu	kaakoo	bata	koa=ra	a	koopu	beara	hogo	va-antee	bata	nao=ra	a	taboo\\
     not.exist	\textsc{nspec}3.\textsc{sg}	measure	because	\textsc{h}-1\textsc{pl}.\textsc{in}.\textsc{pron}	\textsc{tm}2	bend	move	\textsc{cont}	just=1\textsc{pl}.\textsc{in}.\textsc{ipfv}	\textsc{art}3.\textsc{pl}	koopu.bamboo	when=1\textsc{pl}.\textsc{in}.\textsc{pron}	push	\textsc{caus}-stop.at	\textsc{cont}	thither=1\textsc{pl}.\textsc{in}.\textsc{ipfv}	\textsc{art}2.\textsc{sg}	wall\\
\glt `There is not any measure, because we just bend and move the \textit{koopu} bamboo strips along, when we are pushing (them) to the (end of the) wall.'
\z
